% ============================================================================
%  Chapter 14 -- Elementary Machines: Registers and Buses        (Lecture 14)
% ============================================================================
\chapter{Elementary Machines: Registers and Buses}
\label{ch:registers}

\epigraphline{Roads? Where we're going, we don't need roads.}
  {Dr.\ Emmett Brown, \emph{Back to the Future}, Robert Zemeckis, 1985}

\begin{chapgoals}
By the end of this chapter you will be able to
\begin{itemize}
\item describe a register as a set of flip-flops with a common clock, and explain its load signal;
\item explain what a bus is, and compare dedicated buses with a common bus;
\item say why a common bus needs three-state buffers, and what high impedance means;
\item describe a shift register and its four combinations of serial and parallel input and output;
\item convert a datum from parallel to serial and back, in hardware and in \riscv assembly.
\end{itemize}
\end{chapgoals}
% Generated by tools/textbook/export_listings.js from tools/decks/mastery.js and tools/concepts/registry.json: do not edit.
\begin{chapstatus}{Learning status of this chapter}
\masterychip{now}{UNDERSTAND NOW}\enspace shift registers; registers as flip-flops; buses and three-state gates.\par\smallskip
\masterychip{use}{USE NOW}\enspace the complete program, the PRG32 tools (used, not studied, in this course).\par\smallskip
\masterychip{later}{UNDERSTAND LATER}\enspace memory-mapped I/O (Chapter~\ref{ch:memory}); memories and addresses (Chapter~\ref{ch:memory}).
\end{chapstatus}
\begin{chapstatus}{Roadmap of this chapter}
By the end of the chapter you can:
\begin{enumerate}[nosep,leftmargin=2em]
\item explain how registers share a bus through three-state buffers
\item use a shift register for serial transfer
\end{enumerate}
It builds on: registers as flip-flops (Chapter~\ref{ch:sequential}). Next: Chapter~\ref{ch:memory}.
\end{chapstatus}


% ----------------------------------------------------------------------------
\section{Warm-up on the machine}

The display and the sound chip of the kit are connected to the processor by a
few wires, not by thirty-two. A word must therefore travel one bit after
another: a \term{serial} transfer. Listing~\ref{lst:rb:serial} sends a byte
in this way.

\asmfile[caption={Send a byte one bit at a time, most significant bit first.},label=lst:rb:serial]{code/l14/c03.S}

The loop is a shift register in software: shift, take one bit, repeat. (The
two lines left out would write to registers of the chip that control its
pins, a subject of Chapter~\ref{ch:memory}.) The ESP32-C6 contains hardware
that does exactly this, millions of times per second: every pixel of the
display of the kit arrives one bit at a time. This chapter builds the
machines that do it.

% ----------------------------------------------------------------------------
\section{Word gates and registers}

\subsection{Word gates}

A \term{word} is a vector of $n$ bits, $V=(v_0,v_1,\dots,v_{n-1})$. A
\term{word gate} acts on all the bits of a word at once: with a control
signal $a$, it computes $a\cdot v_i$ for every $i$. With $a=1$ the word
passes unchanged; with $a=0$ every bit is forced to 0. It is $n$ identical
gates that share one control input, the building block of every datapath. The
multiplexer on words of Chapter~\ref{ch:combinational} is built from word
gates, and the bitwise instructions \ins{and}, \ins{or} and \ins{xor} of
\riscv are word gates thirty-two bits wide.

\subsection{The register}

A \term{register} is a device able to store a group of bits. In its simplest
form it is a group of flip-flops with a common synchronisation signal, a clock
or an enable (Chapter~\ref{ch:sequential}). \term{Loading} a register means
storing the datum present on its inputs, which happens when the \term{load}
signal is active, typically in step with the clock.

The classical example is the 74273, a register of 8 bits in one chip
(Figure~\ref{fig:rb:chip}): eight D flip-flops, one clock and one clear,
inputs $D_0$ to $D_7$ and outputs $Q_0$ to $Q_7$. In \riscv each of the
registers \reg{x1} to \reg{x31} is thirty-two such flip-flops.

\begin{figure}[htb]
\centering
\begin{subfigure}[b]{0.45\textwidth}\centering
\includegraphics[width=\textwidth,height=36mm,keepaspectratio]{chip_74273.jpg}
\caption{A 74LS273: eight D flip-flops}\label{fig:rb:chip}
\end{subfigure}\hfill
\begin{subfigure}[b]{0.45\textwidth}\centering
\includegraphics[width=\textwidth,height=36mm,keepaspectratio]{pcb.jpg}
\caption{Copper tracks on a printed circuit board}\label{fig:rb:pcb}
\end{subfigure}
\caption{A register and a bus, as objects.}
\end{figure}

% ----------------------------------------------------------------------------
\section{Buses}

\subsection{Data transfers and the bus}

\begin{definition}{Bus}
A \term{bus} is a communication channel that lets the components of an
electronic system exchange data by transmitting and receiving signals: the set
of lines that carries a datum.
\end{definition}

The name is a contraction of the Latin \emph{omnibus}, ``for all''. A bus of
$n$ lines carries a word of $n$ bits in one step. The transfer of a datum
from a register $X$ to a register $Y$ is the fundamental operation of the
architecture of a computer: almost every instruction of
Chapter~\ref{ch:firstcontact} is, at bottom, a sequence of register transfers.

\subsection{Two technical solutions}

With \term{dedicated buses} there is one set of lines for each possible
transfer. Several transfers can happen at the same time; the solution is
efficient but expensive, because the wires grow with the square of the number
of registers. With a \term{common bus} a single path collects many potential
transfers: each time, one source and one destination are chosen, and a
transfer takes two steps, from the source to the bus and from the bus to the
destination. The common bus is cheap and widespread; it is less efficient,
because transfers cannot be simultaneous. Real processors use intermediate
solutions: several buses, each shared by a group of registers.

\subsection{Parallel, serial and mixed transfers}

In a \term{parallel} transfer the $n$ bits of the word travel together on $n$
lines, in one step. In a \term{serial} transfer the bits travel one after
another on a single line, in $n$ steps. A transfer can also be \emph{mixed}:
serial on the byte, for example, with eight lines that carry one byte at a
time.

Buses come in very different physical sizes (Figure~\ref{fig:rb:pcb}). Inside
a chip they are parallel. Between chips and between machines they are today
almost always serial: USB, SATA, PCI Express, HDMI and Ethernet are all
serial, because at very high speeds it is easier to keep one fast line clean
than to keep thirty-two lines in step.

\subsection{From bus to register, and from register to bus}

Figure~\ref{fig:rb:busreg} shows the two elementary operations on a common
bus. \emph{From bus to register}: the datum is loaded into the register $R$
when the load signal $c$ is active (a \term{strobe}, usually in step with the
clock). \emph{From register to bus}: the content of $R$ reaches the bus
through a \term{buffer} controlled by an enable signal; when the enable is not
active the register must leave the bus free for the others. How it does so is
the subject of the next section.

\begin{figure}[htb]
\centering
\begin{tikzpicture}
  \draw[prgamber,line width=3pt] (0,0) -- (5,0) node[pos=0.1,above,lbl,text=prgamber]{bus A};
  \node[blkacc,minimum width=26mm,minimum height=9mm] (r1) at (2.5,-1.6) {register $R$};
  \draw[->,bus] (2.5,0) -- (r1.north);
  \draw[->,thick,prgred] (0,-1.6) node[left,lbl,text=prgred]{load $c$} -- (r1.west);
  \node[lbl] at (2.5,-2.5) {from bus to register};
  \draw[prgamber,line width=3pt] (7,0) -- (12,0) node[pos=0.1,above,lbl,text=prgamber]{bus A};
  \node[blkacc,minimum width=26mm,minimum height=9mm] (r2) at (9.5,-2.2) {register $R$};
  \node[blkvio,minimum width=14mm,minimum height=6mm] (b) at (9.5,-1.0) {buffer};
  \draw[->,bus] (r2.north) -- (b.south);
  \draw[->,bus] (b.north) -- (9.5,0);
  \draw[->,thick,prgred] (7.3,-1.0) node[left,lbl,text=prgred]{enable} -- (b.west);
  \node[lbl] at (9.5,-3.1) {from register to bus};
\end{tikzpicture}
\caption{The two elementary transfers on a common bus.}
\label{fig:rb:busreg}
\end{figure}

% ----------------------------------------------------------------------------
\section{Three-state buffers}

\subsection{The three-state gate}

If two registers drove the same wire, one with 0 and one with 1, the result
would be a short circuit between the supply and ground.

\begin{definition}{Three-state output}
A \term{three-state} (or \emph{tri-state}) output has a third condition
besides 0 and 1: \term{high impedance}, written $z$. In high impedance the
output is electrically disconnected: it neither drives nor loads the wire.
\end{definition}

An enable signal $a$ decides: with $a=1$ the gate works normally, with $a=0$
its output is $z$. Electrically, both output transistors of the gate
(Chapter~\ref{ch:electronics}) are switched off. The table below describes a
three-state AND with an active-high enable.

\begin{center}\small\ttfamily
\begin{tabular}{ccc|c}
\textrm{$a$}&\textrm{$x$}&\textrm{$y$}&\textrm{output}\\\hline
1&0&0&0\\1&0&1&0\\1&1&0&0\\1&1&1&1\\\rowcolor{prgamber!15}0&\textrm{any}&\textrm{any}&z
\end{tabular}
\end{center}

It is the most widespread technique for connecting several source registers
to a common bus.

\subsection{Enabling a bus}

A three-state \emph{word} gate placed between a bus $X$ and a bus $Y$, with an
enable signal $a$, makes $Y$ reproduce $X$ if $a$ is active; otherwise all the
lines of $Y$ are at high impedance, and $Y$ takes the neutral value. One says
that $Y$ is \emph{sustained} by a three-state buffer.

\begin{prgmachine}
The GPIO pins of the ESP32-C6 work in this way: each pin can be an output, or
can be released and used as an input. On the kit the pins of the joystick are
configured as inputs, with their output buffers in high impedance, and the
pins of the display as outputs.
\end{prgmachine}

\subsection{Connecting buses: the bus OR}

Several source registers are connected to one bus simply by joining their
outputs, each through a three-state buffer (Figure~\ref{fig:rb:busor}).

\begin{figure}[htb]
\centering
\begin{tikzpicture}
  \foreach \i/\n/\e in {0/A/\alpha,1/B/\beta,2/X/\xi}{
    \node[blkacc,minimum width=22mm,minimum height=8mm] (r\i) at (\i*3.3,0) {register \n};
    \node[blkvio,minimum width=14mm,minimum height=6mm] (t\i) at (\i*3.3,-1.2) {3-state};
    \draw[->,bus] (r\i) -- (t\i);
    \draw[->,bus] (t\i.south) -- (\i*3.3,-2.2);
    \draw[->,thick,prgred] (\i*3.3-1.5,-1.2) node[left,lbl,text=prgred]{$\e$} -- (t\i.west);
  }
  \draw[prgamber,line width=3pt] (-1.8,-2.2) -- (8.4,-2.2) node[right,lbl,text=prgamber]{common bus $Y$};
  \node[blk,fill=prgsteel,text=white,minimum width=34mm,minimum height=8mm] (d) at (3.3,-3.5) {destination registers};
  \draw[->,bus] (3.3,-2.2) -- (d.north);
\end{tikzpicture}
\caption{Three sources on a common bus, each behind a three-state buffer with
its own enable.}
\label{fig:rb:busor}
\end{figure}

\begin{prgconcept}[The rule of the common bus]
Of the $n$ sources that flow into a bus, at most one is enabled at a time: one
active signal, and the other $n-1$ neutral.
\end{prgconcept}

Inside modern chips long shared lines are usually replaced by multiplexers,
which make the same selection without high-impedance states; between chips,
three-state buses are still the rule.

% ----------------------------------------------------------------------------
\section{Shift registers}

\subsection{Definition}

\begin{definition}{Shift register}
A \term{shift register} is a chain of one-bit memory cells connected to each
other: at every clock pulse each bit moves from one cell to the adjacent one.
\end{definition}

Each cell is a D flip-flop whose input is the output of its neighbour
(Figure~\ref{fig:rb:sipo}). The shift can go in one fixed direction or in a
direction chosen by a control line (bidirectional registers). Edge-triggered
flip-flops are essential here: with transparent latches a bit would run
through the whole chain during one clock pulse.

\begin{figure}[htb]
\centering
\begin{tikzpicture}
  \foreach \i in {0,1,2,3}{
    \node[blkacc,minimum width=17mm,minimum height=11mm] (f\i) at (\i*2.9,0) {D\hfill Q};
    \draw[->,thick] (f\i.south) -- ++(0,-7mm) node[below,font=\ttfamily\small]{Q\i};
    \draw[->,thick,prgred] (\i*2.9-0.5,-2.1) -- (\i*2.9-0.5,-0.55);
  }
  \draw[->,thick] (-2.1,0) node[left,lbl]{serial in} -- (f0.west);
  \foreach \a/\b in {0/1,1/2,2/3} \draw[->,thick] (f\a.east) -- (f\b.west);
  \draw[->,thick] (f3.east) -- ++(1.0,0) node[right,lbl]{serial out};
  \draw[thick,prgred] (-2.1,-2.1) node[left,font=\ttfamily]{CK} -- (3*2.9-0.5,-2.1);
\end{tikzpicture}
\caption{A shift register of four bits built with D flip-flops, with serial
input and both serial and parallel outputs.}
\label{fig:rb:sipo}
\end{figure}

\subsection{Four kinds of shift register}

The input and the output of a shift register can each be serial or parallel,
which gives four kinds (Table~\ref{tab:rb:four}). They make the reading and
the writing of data easier, because they turn parallel data into serial data
and back.

\begin{table}[htb]
\centering\small
\caption{Four kinds of shift register.}\label{tab:rb:four}
\begin{tabularx}{\textwidth}{@{}llL@{}}
\toprule
\textbf{Kind} & \textbf{Name} & \textbf{Use}\\
\midrule
serial in, serial out & SISO & a delay of $n$ clock pulses; a serial link between two registers\\
serial in, parallel out & SIPO & conversion from serial to parallel: the receiver\\
parallel in, serial out & PISO & conversion from parallel to serial: the transmitter\\
parallel in, parallel out & PIPO & an ordinary register that can also shift\\
\bottomrule
\end{tabularx}
\end{table}

\subsection{Transfers between registers}

In a serial transfer the serial output of one shift register feeds the serial
input of another: $n$ clock pulses move a word of $n$ bits, and one wire is
enough. In a parallel transfer the parallel outputs of one register feed the
parallel inputs of another: one clock pulse moves the whole word, and $n$
wires are needed.

\begin{prgconcept}
Moving a word of $n$ bits costs either time or copper: $n$ pulses or $n$
wires. It is the same trade-off as between dedicated and common buses.
\end{prgconcept}

\subsection{Serial to parallel and back}

A complete serial link works in five steps (Figure~\ref{fig:rb:serpar}). The
datum of $n$ bits is placed in parallel in a register $R$. $R$ is a shift
register, and the $n$ bits leave it one after another, in $n$ successive
steps. A single line carries them to the equipment downstream. There another
register is loaded serially in $n$ steps, rebuilding the datum bit by bit.
Finally the rebuilt datum is passed on in parallel.

\begin{figure}[htb]
\centering
\begin{tikzpicture}[node distance=4mm]
  \node[blkacc,text width=21mm,minimum height=15mm] (a) {\textbf{Load}\\[1pt]\footnotesize in parallel};
  \node[blkamb,text width=21mm,minimum height=15mm,right=of a] (b) {\textbf{Shift out}\\[1pt]\footnotesize $n$ steps};
  \node[blkvio,text width=21mm,minimum height=15mm,right=of b] (c) {\textbf{Travel}\\[1pt]\footnotesize one line};
  \node[blkgrn,text width=21mm,minimum height=15mm,right=of c] (d) {\textbf{Shift in}\\[1pt]\footnotesize $n$ steps};
  \node[blkred,text width=21mm,minimum height=15mm,right=of d] (e) {\textbf{Use}\\[1pt]\footnotesize in parallel};
  \foreach \x/\y in {a/b,b/c,c/d,d/e} \draw[->,thick] (\x) -- (\y);
\end{tikzpicture}
\caption{The conversion of a datum from parallel to serial and back.}
\label{fig:rb:serpar}
\end{figure}

This is the hardware of the warm-up. The display of the kit receives each
pixel in this way over a bus called SPI, and the keyboard of
Chapter~\ref{ch:io} sends each key code in this way. UART, SPI, I2C, I2S and
USB all rest on a pair of shift registers, one at each end.

\subsection{Registers with an address: a first memory}

Take four registers of three bits and put a decoder in front of them
(Figure~\ref{fig:rb:mem}). Two address lines, $A_1$ and $A_0$, go to the
2-to-4 decoder, which selects one register. Three data inputs, $I_2I_1I_0$,
reach all the registers, but only the selected one is loaded; three data
outputs, $O_2O_1O_0$, show the content of the selected register through
three-state buffers. Three control lines govern the whole: a \term{chip
select} (CS), which enables the device, a line that chooses between reading
and writing, and an \emph{output enable} (OE). The result is a memory of 4
words by 3 bits.

\begin{figure}[htb]
\centering
\begin{tikzpicture}
  \node[blkamb,minimum width=16mm,minimum height=31mm] (dec) at (0,-1.17) {decoder\\2-to-4};
  \foreach \i in {0,1,2,3}{
    \node[blkacc,minimum width=36mm,minimum height=6.5mm] (r\i) at (3.7,-\i*0.78) {register \i\ (3 bits)};
    \draw[->,thick,prgamber] (dec.east|-r\i.west) -- (r\i.west);
  }
  \draw[->,thick] (-2.2,-1.17) node[left,font=\ttfamily]{A1 A0} -- (dec.west);
  \draw[->,bus] (3.7,1.0) node[above,font=\ttfamily]{I2 I1 I0} -- (r0.north);
  \draw[->,bus] (r3.south) -- (3.7,-3.4) node[below,font=\ttfamily]{O2 O1 O0};
  \node[lbl,text=prgred] at (-0.6,-3.4) {control: CS, RD, OE};
\end{tikzpicture}
\caption{Four registers and a decoder: a memory of 4 words by 3 bits.}
\label{fig:rb:mem}
\end{figure}

\begin{prgconcept}
A decoder in front of a group of registers makes a memory. The next chapter
scales this up to the main memory of a computer.
\end{prgconcept}

% ----------------------------------------------------------------------------
\section{Back to the machine}

\subsection{Buses on the PRG32 kit}

The kit contains an example of nearly every kind of bus of this chapter
(Table~\ref{tab:rb:kit}). The chip select of the display is the CS of the
small memory just seen: several devices can share the same serial bus because
only the selected one listens.

\begin{table}[htb]
\centering\small
\caption{Buses on the \prg kit (from \fname{docs/hardware} of the
repository).}\label{tab:rb:kit}
\begin{tabularx}{\textwidth}{@{}P{0.21\textwidth}P{0.16\textwidth}LL@{}}
\toprule
\textbf{Connection} & \textbf{Kind} & \textbf{Lines} & \textbf{What travels on it}\\
\midrule
inside the ESP32-C6 & parallel & 32 data lines, address and control & instructions, data and register transfers\\
display (ILI9341) & serial, SPI & clock, data out, chip select, data/command & every pixel of every frame, bit by bit\\
sound amplifier & serial, I2S & bit clock, word select, data & audio samples, 22\,050 per second\\
joystick and buttons & one wire per button & one input pin each & a decoded representation: one line per button\\
USB to the laptop & serial, differential pair & two data wires & the firmware and the debug console\\
\bottomrule
\end{tabularx}
\end{table}

\subsection{Shifts in RISC-V: a shift register in one instruction}

\ins{slli} and \ins{srli} move all the bits of a register by a given number
of positions, filling with zeros; \ins{srai} fills with copies of the sign
bit (Chapter~\ref{ch:integers}); \ins{sll}, \ins{srl} and \ins{sra} take the
amount from a second register. With them the two conversions are two
instructions each (Listing~\ref{lst:rb:shift}).

\asmfile[caption={Parallel to serial, and serial to parallel.},label=lst:rb:shift]{code/l14/c05.S}

A true shift register needs $n$ clock pulses to shift by $n$ positions. The
processor does not work in that way: in hardware a shift by any amount in one
step is done by a \term{barrel shifter}, which is made of layers of
multiplexers, not of a chain of flip-flops. It is combinational, it does its
work within one clock period, and it is one of the units of the ALU of
Chapter~\ref{ch:combinational}.

% ----------------------------------------------------------------------------
% Generated by tools/textbook/export_listings.js from decks/plan.js and code/l14.S: do not edit.
\section{The program of this chapter}\label{sec:lab:l14}

The listings of this chapter are parts of one program, \fname{code/l14.S} in the
repository of the course, printed in full in Listing~\ref{lst:lab:l14}. The
automatic tests of the course material build it and run it on the headless
runner of \prgqt{}, with scripted input, and check the frame;
Figure~\ref{fig:lab:l14} is the frame that the test wrote.

\begin{description}[style=nextline,leftmargin=1.5em]
\item[On the screen.] Top row: the eight bits of 0xA5 in the order in which they leave, most significant first (lit, dark, lit, dark, dark, lit, dark, lit). Bottom row: the byte rebuilt by the receiver. A green square says that it is equal to the byte sent.
\item[What to press.] Nothing: the program runs by itself.
\end{description}

\begin{figure}[htb]
\centering
\includegraphics[width=0.5\textwidth]{lab_l14.png}
\caption{\fname{l14.S} on \prgqt{}.}\label{fig:lab:l14}
\end{figure}

\asmfile[caption={The complete program of the chapter, \fname{code/l14.S}.},label=lst:lab:l14]{code/l14/program.S}

To run it on \prgqt{} the commands are those of Listing~\ref{lst:lab:l14:run};
Section~\ref{sec:prg:testing} gives the same steps for QEMU and for the board.

\begin{shell}[caption={Building \fname{l14.S} and running it on \prgqt{}.},label=lst:lab:l14:run]
# 1. start PRG32-QT on this computer (it listens on port 8080)

# 2. in the PRG32 folder, with ESP-IDF active: build the cartridge
python3 -m prg32 cartridge build ~/arc12/code/l14.S --portable \
    --entry-prefix l14 --name l14 --out l14.prg32

# 3. send it to PRG32-QT and start it
python3 -m prg32 esp32c6 upload-and-run l14.prg32 \
    --url http://localhost:8080 --slot cart0

# 4. look, press, compare with the previous slide
\end{shell}


% ----------------------------------------------------------------------------
% Generated by tools/textbook/export_listings.js from tools/concepts/assessments.json: do not edit.
\begin{chapstatus}{The thread of this chapter}
A bus is a road that only one register may drive at a time, and a shift register turns a word into a stream of bits: the display of the kit receives its pixels this way. Next: memory.
\par\medskip\textbf{Exit ticket} (two to four minutes, alone, on paper):
\begin{enumerate}[nosep,leftmargin=2em]
\item Three registers share a common bus. What happens if two of them drive it at the same time, and which component prevents it?
\item An 8-bit word is sent serially, one bit per clock pulse. How many pulses does it take, and what does the shift register do at each?
\end{enumerate}
\end{chapstatus}


\section{Check your understanding}

\begin{quiz}{Why does a common bus need three-state buffers?}
\item To make transfers faster
\item So that only the enabled source drives the lines
\item To store the datum on the bus
\item To convert serial data to parallel
\end{quiz}

\begin{quiz}{How many clock pulses does a serial transfer of a 32-bit word need?}
\item 1
\item 5
\item 32
\item 64
\end{quiz}

\begin{quiz}{A memory has 10 address lines and 8 data lines. How many bits does it store?}
\item 80
\item 1\,024
\item 8\,192
\item 10\,240
\end{quiz}

\section{Exercises}

\begin{exercise}[paper first, then \prg]
\begin{enumerate}[nosep]
\item Draw a 4-bit serial-in, parallel-out shift register with D flip-flops,
  and fill in its contents over four clock pulses for the input 1, 0, 1, 1.
\item Three registers share a common bus: write the enable and load signals
  needed to copy register $B$ into register $A$.
\item Complete the warm-up: write the loop that receives eight bits from a
  line and rebuilds the byte.
\item On \prg, show a byte as eight lamps and shift it to the left by one
  position every half second.
\item Make the bit that leaves on the left come back in on the right (a
  rotation): a \emph{ring counter}.
\end{enumerate}
\expected{The third task is \code{slli a1, a1, 1} and then \ins{or} with the
received bit, eight times. The last is \code{srli t1, a0, 7};
\code{andi t1, t1, 1}; \code{slli a0, a0, 1}; \code{or a0, a0, t1};
\code{andi a0, a0, 255}. On \prg: a pattern of lamps that walks to the left
and, in the last task, goes round for ever.}
\end{exercise}

\section{Glossary}

\begin{glossarytable}
word gate & porta di parola\\
register & registro\\
bus & bus\\
dedicated / common bus & bus dedicato / comune\\
load, strobe & caricamento, strobe\\
three-state (tri-state) & a tre stati (tristate)\\
high impedance & alta impedenza\\
shift register & registro a scorrimento\\
serial / parallel transfer & trasferimento seriale / parallelo\\
chip select & selezione del chip\\
\end{glossarytable}

\section{Conclusions and reading}

Registers are an evolved form of memory with an elementary way of working.
Reading and writing them happens through buses of different kinds. For a
common bus to work, three-state circuits are needed: 0, 1, and disconnected.
Data can be transferred serially or in parallel, and shift registers convert
between the two. A decoder in front of a group of registers makes a memory:
the subject of the next chapter.

\paragraph{Reading.} \citet{patterson2020}, appendix~A; \citet{bucci2017},
the chapters on registers and buses; \fname{docs/hardware/hardware.md} in the
\prg repository.
